\chapter*{Preface}
\addcontentsline{toc}{chapter}{Preface}

Mathematics serves as the quintessential language of physics. From classical Newtonian trajectories to general relativity and quantum field theories, physical laws achieve precision, predictive power, and beauty through mathematical formulation. 

This textbook, \textbf{Mathematics for Physics}, is designed specifically for undergraduate and beginning graduate students in physics, applied physics, and related engineering disciplines. The core mission of this book is to bridge the gap between abstract mathematical formalism and concrete physical intuition. Rather than treating mathematical theorems in isolation, every technique introduced herein is directly motivated by real-world physical systems---including electromagnetic field distributions, planetary motion, oscillating circuits, and quantum state evolutions.

\section*{Pedagogical Framework: The Active Learning Cycle}

To maximize conceptual retention and analytical proficiency, each chapter integrates a structured four-stage pedagogical cycle:
\begin{enumerate}
    \item \textbf{Conceptual Motivation:} Clear physical context establishing why a specific mathematical tool is necessary.
    \item \textbf{Rigorous Analytical Development:} Thorough derivations that maintain mathematical integrity without sacrificing accessibility.
    \item \textbf{P-S-C-T Worked Examples:} Problem-solving modeled through four explicit steps:
    \begin{itemize}
        \item \textbf{Picture (P):} Physical conceptualization, coordinate selection, and free-body/field sketches.
        \item \textbf{Strategy (S):} Identifying applicable physical principles, conservation laws, or symmetry arguments.
        \item \textbf{Calculation (C):} Step-by-step mathematical execution with dimensional verification.
        \item \textbf{Think (T):} Physical interpretation of limiting cases, asymptotic behaviors, and practical implications.
    \end{itemize}
    \item \textbf{Numerical Integration \& Python Modeling:} Modern scientific computing scripts using NumPy, SciPy, and Matplotlib to visualize complex dynamics and solve differential equations beyond closed-form limits.
\end{enumerate}

\section*{How to Use This Book}

Chapters 1 through 2 establish the geometric and differential foundations of curvilinear coordinates and vector calculus. Chapters 3 and 4 explore infinite series expansions and complex analysis, vital for optics and wave mechanics. Chapter 5 develops matrix mechanics and eigenvalue techniques. Chapters 6 and 7 provide a rigorous treatment of multivariable calculus and the fundamental integral transformation theorems of Gauss, Stokes, and Green. Finally, Chapters 8 and 9 focus on first- and second-order ordinary differential equations, emphasizing damped, driven, and resonant physical oscillators.

I express my deepest appreciation to my colleagues in the Department of Physics at Rambhai Barni Rajabhat University, whose scholarly feedback has greatly enriched this work, and to the dedicated physics students whose insightful questions continue to shape its evolution.

\vspace{1.5cm}
\begin{flushright}
    \textbf{Asst. Prof. Dr. Chewa Thassana}\\
    Department of Physics, Faculty of Science and Technology\\
    Rambhai Barni Rajabhat University (RBRU)\\
    Chanthaburi, Thailand
\end{flushright}
